\subsection{Operational Reproducibility and Correctness}
\label{sec:operational-reproducibility}

Deciding whether a model can be relied on for a specific operational decision requires
a statement about that decision, not about a distribution of them. Aggregate accuracy
does not provide one: a model that is always correct on one input and always wrong on
another scores the same as a model that is inconsistent on both, although the two fail
in different ways and only one of them fails visibly. We therefore define a criterion
that is evaluated where the decision is made.

We define reproducibility over a test matrix of models, inputs, and
interventions declared before evaluation, and assess it within each condition
over its $R$ executions rather than averaging across the matrix. Correctness is
fixed by task ground truth established independently of model responses: the
car-wash objective requires transporting the car, so \texttt{DRIVE} is correct.
We use $R=10$ single-pass sampled executions per condition at temperature $1.0$,
since greedy decoding repeats one action and conceals behavioral variation. Two
models do not accept a temperature parameter and ran at provider default. Invalid
or missing outputs never satisfy the criterion, which applies to the declared
matrix alone.

Evaluating within conditions rather than across them is not stricter bookkeeping.
An accuracy figure aggregated over a matrix is a property of the inputs chosen for
that matrix as much as of the model: in the results below, the same question under
different surrounding text produces a correct-response rate ranging from $7.6\%$ to
$95.9\%$. Reported per condition, each figure describes a declared input;
aggregated, it describes a set of inputs that someone selected.

Let $y_r$ denote the operational outcome of execution $r$ under
a fixed condition, and $y^*$ the predefined correct action.
Reproducible correctness requires
\begin{equation}
    y_r = y^* \quad \forall r \in \{1,\ldots,R\}.
    \label{eq:reproducible-correctness}
\end{equation}

Ten identical incorrect decisions are reproducible but wrong;
mixed decisions are non-reproducible. Aggregate accuracy cannot
distinguish these states and is therefore reported separately
from the number of conditions meeting (\ref{eq:reproducible-correctness}).

\begin{table}[t]
\centering
\caption{Classification of outcome patterns for the car-wash task
($R=10$; correct action: DRIVE).}
\label{tab:operational-states}
\vspace{6pt}

\small
\begin{tabular*}{0.75\columnwidth}{@{\extracolsep{\fill}}ll@{}}
\toprule
\textbf{Observed outcomes} & \textbf{Reproducibility} \\
\midrule
10 DRIVE / 0 WALK & Reproducibly correct \\
10 WALK / 0 DRIVE & Reproducibly incorrect \\
7 WALK / 3 DRIVE & Non-reproducible \\
\bottomrule
\end{tabular*}
\end{table}

\subsection{Semantic Interventions and Output Preference}
\label{sec:semantic-interventions}

We study a fixed operational task under different natural-language
interventions. All modify the model's input without changing parameters,
but differ in what they specify: conventional prompts request reasoning
procedures, expertise, or emphasis, whereas the substrate specifies the
relations among the task objective, the required object, and the available
actions.

Models are evaluated behaviorally: for every model and condition, we record
the emitted output of each execution and classify it to an operational
action. Two additional measurements are available on subsets. Where token
log-probabilities are exposed, including for closed-weight models, we
measure relative preference between the two actions; for open-weight
models, we additionally access internal activations. Model identifiers,
access levels, inference configurations, and exact intervention texts
are given in Appendix~\ref{app:experimental-details}.

Where log-probabilities are available, we measure relative preference
at the position predicting the first answer token:
\begin{equation}
    M = \log \sum_{t \in \mathcal{D}} P(t)
      - \log \sum_{t \in \mathcal{W}} P(t),
    \label{eq:decision-margin}
\end{equation}
where $\mathcal{D}$ and $\mathcal{W}$ contain the distinct token IDs
corresponding to the measured surface forms of each action, including
case and leading-space variants. Positive $M$ favors DRIVE. An
intervention's effect is measured relative to the untreated task under
the same model and configuration:
\begin{equation}
    \Delta M(u) = M(u) - M(\emptyset).
    \label{eq:intervention-effect}
\end{equation}

When a provider exposes only the top $K$ tokens, incomplete coverage
of an action's surface forms prevents exact margin estimation. We
report a one-sided bound only when its direction is established by
the available probabilities; otherwise, the margin is recorded as
unavailable. Behavioral correctness is measured over all $R$ declared
executions, with invalid or missing outputs counted as failures and
reported separately.

The emitted action, not the margin, is the criterion. A margin can move
substantially while the emitted action does not: several of the interventions
reported below shift $M$ by nats on models whose emitted word is unchanged, and in
a small number of conditions the sign of $M$ and the emitted token disagree
outright. Treating the margin as the outcome would record those as effects. We
therefore report $M$ as evidence about relative output preference and classify
conditions by the emitted action, which does not depend on log-probability access.

\subsection{Experimental Design}
\label{sec:experimental-design}

The study proceeds in two stages. First, we evaluate the untreated
task, conventional semantic interventions, and a shared task-specific
specification across the declared model matrix. Each condition is
evaluated through repeated single-pass executions, without retries,
majority voting, or post-generation correction. We measure the
resulting operational states and determine which model--condition
pairs satisfy the reproducible-correctness criterion of
\S\ref{sec:operational-reproducibility}.

The order of the two stages is a requirement of the second, not a convenience.
Comparing internal states across models presumes that the models are doing the same
thing, and under an untreated task they are not: which models appear anomalous
depends on the task chosen and on which models were tested, and the difference
between a model that selects the intended action on three of ten executions and one
that selects it on none has no fixed standard against which either is measured. The
alternative usually available --- restricting attention to cases on which the models
already agree --- substitutes a selection bias for the confound, since the comparison
is then limited to whatever the models happen to find easy. Establishing behavioral
convergence first supplies the standard. It is verified per model--condition before
any activation is inspected, and it holds the input, the emitted action, and the
correctness criterion fixed across the compared models by construction, so remaining
internal differences are attributable to the models rather than to divergent task
behavior.

Second, we investigate the engineered behavioral transition internally
in open-weight models. Holding the model, task, and inference
configuration fixed, we compare the untreated and substrate-conditioned
computations. At selected layers, we transplant the substrate-conditioned
activation at the final prediction position into the untreated
computation and measure its effect on the final decision margin.
This tests where the induced internal state becomes sufficient,
under the remaining computation, to change the model's output
preference. Only model--conditions in which the specification changes the emitted
action enter this analysis; a model that does not converge under the specification
is not a missing data point but a separate outcome of the same procedure, and is
reported as one.
